%%%%%%%%%%%%%%%%%%%%%%%%%%%%%%%%%%%%%%%%%%
% Medium Length Professional CV
% LaTeX Template
% Version 2.0 (8/5/13)
%
% This template has been downloaded from:
% http://www.LaTeXTemplates.com
%%
% Original author
% Trey Hunner (http://www.treyhunner.com/)
%
% Important note:
% This template requires the resume.cls file to be in the same directory as the
%% .tex file. The resume.cls file provides the resume style used for structuring thez
%% document.

%
%%%%%%%%%%%%%%%%%%%%%%%%%%%%%%%%%%%%%%%%%

%----------------------------------------------------------------------------------------
%    PACKAGES AND OTHER DOCUMENT CONFIGURATIONS
%----------------------------------------------------------------------------------------


\documentclass{resume} % Use the custom resume.cls style

\name{Zhenzhong (Jack) Xing} % Your name
\address{Ithaca, New York 14850} % Your address
%\address{} % Your secondary addess (optional)
\address{(641)$\cdot$275$\cdot$6452 \\ zx296@cornell.edu \\ https://technium7.github.io/} % Your phone number and email





\begin{document}

%----------------------------------------------------------------------------------------
%    EDUCATION SECTION
%----------------------------------------------------------------------------------------

\begin{rSection}{Education}
\item {\bf Cornell University} 

Ph.D. in Electrical and Computer Engineering \hfill {Expected summer of 2027}

M.S. in Electrical and Computer Engineering\hfill {Dec 2025     \em GPA 4.067}

        with concentrations in Applied Physics and Theory of Computation

\item {\bf Grinnell College}

B.A. in Physics (Honors) and Mathematics\hfill {May 2022     \em GPA 4.0}\\ 
\end{rSection}

%----------------------hm------------------------------------------------------------------
%    WORK EXPERIENCE SECTION
%----------------------------------------------------------------------------------------
\begin{rSection}{Experience and Skills}
\begin{tabular}{ @{} >{\bfseries}l @{\hspace{6ex}} l }
Programming & Python, C, C++, Fortran, Shell\\
Tools & Git, PyTorch, JAX, OpenCV, Numpy, Scipy, Tidy3D, GDS factory\\
Software & Fusion360, Altium, MATLAB, Mathematica, Linux\\
Hardware & Precision optics, Analog/digital electronics, Micro-controller, Raspberry pi, Machining
\end{tabular}

\end{rSection}

\begin{rSection}{Selected Courses}

\textbf{Optoelectronics and Physics}: Nonlinear and quantum optics, Photonics, Electrodynamics, Semiconductor physics, Quantum mechanics, Atomic physics, Quantum information theory

\textbf{Electrical engineering and CS}: Digital systems, Embedded systems, Machine learning, Robot learning and perception, Matrix computation, Analysis of algorithms, Optimization and control

\textbf{Mathematics}: Linear and abstract algebra, Real and complex analysis, Probability and statistics


\end{rSection}


\begin{rSection}{Experience}

\textbf{Cornell University}\\
\begin{rSubsection}{Photonics and Quantum Electronics Group}{Jun 2022 - Present}{Graduate research assistant  \textsc{(founding lab member)}}{Ithaca, NY}

\item \textbf{Designed and built a cryogenic, UHV apparatus, a electrical control systems, and a UV-IR precision optical system from the ground up} as the platform for trapping and controlling calcium ions for quantum information science. The system consists of \textbf{multi-stage vacuum chambers, DC and RF filtering and stabilization circuit boards, and a combination of free-space optics, fibers, and integrated photonics} to address the ion with narrow-linewidth lasers of wavelength from 375 nm to 866 nm and end power of mW level, with the electrical and optical control system developed based on FPGA, DDS, and DAC.

\vspace{0.1in}

\item \textbf{Developed and deployed python packages to accelerate the full experiment cycle}, some are currently used in other research groups at ETH and Duke.
\begin{itemize}
    \item[.] Streamlined the design, simulation, and post-processing process of the ion trap chip with user friendly graphic interface, by interacting with Klayout, GDS factory, Nullspace and Tidy3D.
    \item[.] Automated the calibration and long data taking with in situ data analysis of the trapped ion system.
\end{itemize}

\vspace{0.1in}

\item \textbf{Built and trained a table-top 6-axis robotic arm to facilitate trap packaging and optical path alignment}. Based on an open source project, further developed hardware, MCU firmware, and software for robot perception, control, and motion.
\begin{itemize}
    \item[.] Implemented a Python robot control stack through ASCII serial protocol and keyboard teleop with PyBullet IK and simulation to hardware joint mapping for demo collection and debugging.
    \item[.] Developed multi-camera vision pipelines that interact with Realsense, Thorlab, Infinity cameras and utilize calibration with environmental features to achieve sub-cm precision.
    \item[.] Trained policies based on visual detection + collision-aware planning and vision encoder + diffusion for subtasks such as grabbing/inserting optical sensor and epoxy application, with global planning and interface through local llama model.
\end{itemize}

\vspace{0.1in}

\item \textbf{Demonstrated novel laser cooling schemes of trapped ion} that could reduces the total cooling time by $>10\times$ and address the bottleneck in current industrial trapped ion quantum computing systems. Results lead to two first author papers.

\item \textbf{Demonstrated trapped ion transportation on chips with multi-layer integrated photonics for the first time}. Preparing manuscript for publications.

\begin{itemize}
    \item[.] Built a dynamic multi-voltage control solver with quadratic programming. Reduced the complexity from $\mathcal{O}(nm)$ to $\mathcal{O}(n)$ for ion transport and achieved smoother boundary solutions by reformulating and preconditioning the ion transport problem
    \item[.] Applied CEM to optimize transport path that reduce the final simulated phonon number by $\sim 10 \times$.
    
\end{itemize}


\vspace{0.1in}


\end{rSubsection}

\textbf{Fermilab}\\
\begin{rSubsection}{Detector R\&D Gorup in Neutrino Division}{Summer 2021}{Research intern}{Batavia, IL}
\item Optimized the performance of gaseous argon time projection chamber for neutrino detection using Cython with results published in internal reports
\end{rSubsection}

\textbf{Grinnell College}\\
\begin{rSubsection}{Gravitational Physics Group}{Summer 2020}{Research assistant}{Grinnell, IA}
\item Developed formalism for describing quantum particle constrained on 2D manifolds using differential geometry. Results published in \textit{Physics Scripta}.
\end{rSubsection}




%------------------------------------------------

%\begin{rSubsection}{Florida State University}{Summer 2017 (10 weeks)}{Research assistant at the aero-propulsion, \\ mechatronics and energy (AME) research labs.}{Tallahassee, FL}
%\item Developed a code that simulates the surface temperature of complex geometries based on real weather data in MATLAB and Fortran. 
%Advised by Prof. Neda Yaghoobian.
%\end{rSubsection}
%------------------------------------------------





\end{rSection}

%----------------------------------------------------------------------------------------
%    TECHNICAL STRENGTHS SECTION
%----------------------------------------------------------------------------------------



%\newpage

\begin{rSection}{Publications}
\item \textbf{Z. Xing}, H. M. Rivy, V. Natarajan1, A. K. Kolhatkar, G. Beck, and K. K. Mehta. ``Rapid multi-mode trapped-ion laser cooling in a phase-stable standing wave." \textit{Physical Review X} (in review).

\item A. M. Kolhatkar, V. Natarajan, H. M. Rivy, \textbf{Z. Xing}, G. Beck, K. K. Mehta. ``Enhanced trapped-ion quantum control in foundry-fabricated integrated devices." \textit{Quantum Information Science, Sensing, and Computation XVIII}, Jul 2026.

\item \textbf{Z. Xing}, K. K. Mehta.``Trapped-Ion Laser Cooling in Structured Light Fields." \textit{Physical Review Applied}, July 2025.

\item S. Rodriguez, L. Rodriguez, \textbf{Z. Xing}, C. McMillin, L. R. Ram-Mohan.``Geometrodynamics of a 2D Curved Surface Due to a Constrained Quantum Particle via Its Gravitational Dual: S2 Analytical Model Calculations." \textit{Physica Scripta}, May 2025.

%\item \textbf{Oscar Jaramillo}, Leonardo Massai and Karan Mehta. HfO$_2$-based Platform for High Index Contrast
%Visible and UV Integrated Photonics. \textit{In 2023 Conference on Lasers and Electro-Optics (CLEO).}
\end{rSection}








\begin{rSection}{Selected Awards}
Engineering School Fellowship (Cornell University).\\
Andrew W. Archibald Prize for Highest Scholarship (Grinnell College).\\
H. George Apostle Prize in Physics (Grinnell College).\\
\end{rSection}


\begin{rSection}{Languages}

\begin{tabular}{ @{} >{\bfseries}l @{\hspace{6ex}} l }
English & (Proficient) \\
Mandarin & (Proficient)
\end{tabular}
\end{rSection}

%----------------------------------------------------------------------------------------
%    EXAMPLE SECTION
%----------------------------------------------------------------------------------------

%\begin{rSection}{Section Name}

%Section content\ldots

%\end{rSection}

%----------------------------------------------------------------------------------------

\end{document}